\metatitle{Reductions of the Lorentzian conjecture}
\metadescription{Exact binary encodings, reduction of positive weights to
unweighted completion counts, and an equivalent rank-two counting conjecture.}

\section[completionReductions]{Reductions of the Lorentzian conjecture}

We reduce the \href{rank.htm#lorentzianConjecture}{Lorentzian conjecture}
to unweighted counting. It suffices to consider colored matroids in
which each position allows at most two colors and the distinguished
color occurs twice in every word.

Let $\mathcal F$ be a \href{intro.htm#coloredMatroids}{colored matroid}
on $E$, and fix a color $c$. As in the introduction, $B(u,c)$ denotes
the positions carrying color $c$. Its weighted completion polynomial is
\[
 \begin{gathered}
 Q_c(\mathbf x;\mathbf w)=\\
 \sum_{u\in\mathcal F}
 \left(\prod_{e:u_e\ne c}w_{e,u_e}\right)
 \prod_{e:u_e=c}x_e.
 \end{gathered}
\]
The weights $w_{e,a}$ are positive, and the degree is the quota $r_c$.
When all weights are $1$, the coefficient of $\prod_{e\in B}x_e$
counts the words with $B(u,c)=B$.

\section[binaryLists]{Binary lists and exact encodings}

The \defin{color list} at a position is
\[
 L_e\coloneqq\{u_e:u\in\mathcal F\}.
\]
We call $\mathcal F$ a \defin{binary-list family} if $|L_e|\leq2$ at
every position. Such a family may use more than two colors in total.

Delete positions with a single possible color. The remaining
\defin{color graph} has the colors as vertices and each position $e$
as an edge joining the two colors in $L_e$. Orient $e$ toward $u_e$.
Quotas prescribe indegrees, and a \href{intro.htm#cyclicExchange}{cyclic
exchange} reverses a directed simple cycle in both words. We retain
exactly the orientations corresponding to words of $\mathcal F$.

We now replace each position by several positions with binary lists.
Choose positive integers $m_{e,a}$ for $a\in L_e$. Let $E_e$ consist
of the positions $(e,a,j)$ with $a\in L_e$ and $1\leq j\leq m_{e,a}$,
and put $M_e=|E_e|$.
Introduce distinct new colors $p_e$, one for each $e$. To expand a word
$u$, choose one position $(e,u_e,j)$ in each $E_e$, give it color $u_e$,
and give every other position of $E_e$ color $p_e$.
Let $\widehat{\mathcal F}$ consist of all expansions of words in
$\mathcal F$.

\begin{proposition}[Expansion by private colors]
For any nonempty fixed-quota family $\mathcal F$, its expansion is
binary-list and satisfies symmetric cycle exchange if and only if
$\mathcal F$ does. Original colors retain their quotas, and $p_e$ has
quota $M_e-1$ with projection
$U_{M_e-1,M_e}$ on $E_e$ and loops elsewhere. Colors of quota zero may
be omitted. The color graph is bipartite between original and new colors.
If every $m_{e,a}=1$, expansion is a bijection.
\end{proposition}

\begin{proof}
The unique entry different from $p_e$ recovers $u_e$. Every position
of $E_e$ can carry this entry, giving the quotas and uniform projection.

Take two expanded words and prescribe a disagreement in $E_e$.
If the original entries agree, say both equal $a$, the two chosen
positions form the two-cycle $a\to p_e\to a$. Switching them leaves
the original words fixed.

If the original entries are $a\ne b$, their disagreement in $E_e$ is
\[
 a\longrightarrow p_e\longrightarrow b.
\]
Expand each arrow of an original cyclic exchange through $e$ into such
a path. All colors on the expanded cycle are distinct. Switching it
gives expansions of the two resulting original words and includes
either prescribed disagreement in $E_e$.

For the converse, choose one fixed representative $(e,a,1)$ for every
original assignment. Lift two original words using these representatives.
Only differing positions contribute paths through $p_e$. A cycle meeting
$p_e$ must use both arrows of that path. Contracting these paths gives
an original cyclic exchange with both resulting words in $\mathcal F$.
\end{proof}

\subsection[slotCounting]{Counting the expansions}

Distinguish an original color $c$. Give its new positions variables
$x_{e,c,j}$. The unweighted completion polynomial of the expanded family
is
\[
 \sum_{u\in\mathcal F}
 \left(\prod_{e:u_e\ne c}m_{e,u_e}\right)
 \prod_{e:u_e=c}
 \left(\sum_{j=1}^{m_{e,c}}x_{e,c,j}\right).
\]
For each fixed word, the factors count the choices at its positions;
the color-$c$ factors also record their variables.

The construction retains membership in $\mathcal F$, so its image may
omit quota-compatible orientations. The
\href{series-parallel.htm#listStability}{allowed-color stability theorem}
applies when all words with the given lists and quotas are included.

\section[weightsToCounts]{Replacing weights by completion counts}

\begin{theorem}[Weighted and unweighted forms]
The following universal assertions are equivalent.
\begin{enumerate}
\item Every colored matroid has Lorentzian completion polynomials
for every distinguished color and every choice of positive weights.
\item Every colored matroid has Lorentzian unweighted completion
polynomials for every distinguished color.
\item Every binary-list colored matroid with bipartite color graph
has Lorentzian unweighted completion polynomials for every color on
one side, when each color on the other side has projection
$U_{d-1,d}$ on its $d$ possible positions.
\end{enumerate}
\end{theorem}

\begin{proof}
The first two implications are immediate. For the converse, we realize
weights as numbers of expansions.

First suppose that all weights are positive integers. In the expansion
above, take $m_{e,c}=1$ and $m_{e,a}=w_{e,a}$ for $a\ne c$.
Identify the variable of $(e,c,1)$ with $x_e$. The displayed counting
formula is then exactly $Q_c(\mathbf x;\mathbf w)$.

For positive rational weights, choose a positive integer $d_e$ clearing
the denominators at position $e$. Set
\[
 m_{e,c}=1,\qquad m_{e,a}=d_e w_{e,a}\quad(a\ne c).
\]
At positions allowing only $c$, take $d_e=1$. If $G$ is the unweighted
expanded completion polynomial and $D=\prod_{e\in E}d_e$, then
\[
 G(\mathbf x)=D\,Q_c((x_e/d_e)_{e\in E};\mathbf w).
\]
A word with color-$c$ set $B$ acquires the factor
$D/\prod_{e\in B}d_e$. The third assertion applies to the expansion. Positive
scalar multiplication and invertible positive diagonal substitutions
preserve Lorentzianity in both directions.

Finally, approximate positive real weights by positive rational weights.
The coefficients converge, and the Lorentzian cone is closed. These
closure and substitution properties follow from
\cite[Theorems~2.10 and 2.25]{BrandenHuh2020}.
\end{proof}

\begin{example}[Seventeen unweighted words]
Start with the six words on four positions having quotas $(2,2)$.
Distinguish color $1$, and give the color-$2$ assignments weights
$(2,3,1,1)$. Their completion polynomial is
\[
 \begin{aligned}
 Q_1(\mathbf x;\mathbf w)
 ={}&x_1x_2+3x_1x_3+3x_1x_4\\
 &+2x_2x_3+2x_2x_4+6x_3x_4.
 \end{aligned}
\]
Use one color-$1$ position and respectively $2,3,1,1$ color-$2$
positions in the four expanded sets. Label their private colors
$3,4,5,6$. The expansion has eleven positions, quotas
$(2,2,2,3,1,1)$, and seventeen distinct words.
For instance, the color-$1$ set $\{3,4\}$ has six completions: two
choices for color $2$ in $E_1$ and three choices in $E_2$.
\end{example}

\section[quadraticCounting]{An equivalent rank-two counting conjecture}

Suppose $r_c\geq2$. A nonzero derivative of order $r_c-2$ comes from
\href{operations.htm#deletionAndContraction}{pinning color $c$ at
$r_c-2$ positions and deleting them}. The remaining family is a colored
matroid of quota two for $c$. Repeated differentiation in one variable
gives zero because $Q_c$ is multiaffine.

By the \href{intro.htm#coloredMatroids}{color-projection proposition},
the support is matroidal. The quadratic-derivative characterization
\cite[Theorem~2.25]{BrandenHuh2020} therefore reduces the Lorentzian
conjecture to quota two. Degrees zero and one satisfy the definition.

\begin{conjecture}[Rank-two completion counts]
Let $\mathcal F$ be a binary-list colored matroid with a distinguished
color $c$ of quota two. For any $S\subseteq E$ and $j=0,1,2$, put
\[
 N_j\coloneqq\#\{u\in\mathcal F:|B(u,c)\cap S|=j\}.
\]
Then $N_1^2\geq4N_0N_2$.
\end{conjecture}

\begin{theorem}[Counting equivalence]
The rank-two counting conjecture is equivalent to the weighted
Lorentzian conjecture. It suffices to require the counting inequality
for bipartite binary-list families with the distinguished color on
one side and corank-one projections on the other side.
\end{theorem}

\begin{proof}
Lorentzianity implies the counting inequality by substituting $x_e=s$
on $S$ and $x_e=t$ elsewhere. The resulting polynomial is
$N_2s^2+N_1st+N_0t^2$; its Hessian has at most one positive eigenvalue
exactly when $N_1^2\geq4N_0N_2$.

Conversely, let $Q_c$ have degree two and positive rational weights.
Choose positive rational vectors $\mathbf a,\mathbf b$. At position $e$,
choose an integer $d_e$ so that $d_ea_e$, $d_eb_e$, and all $d_ew_{e,h}$
are integers. Where $c\in L_e$, use $d_e(a_e+b_e)$ positions designated
for color $c$, marking $d_ea_e$ of them. For every other allowed color
$h$, use $d_ew_{e,h}$ positions.

The marked unweighted enumerator of this family is
\[
 D\,Q_c(s\mathbf a+t\mathbf b;\mathbf w),
 \qquad D=\prod_{e\in E}d_e.
\]
Write $Q_c(\mathbf x;\mathbf w)=\tfrac12\mathbf x^{\mathsf T}H\mathbf x$.
The counting inequality gives
\[
 (\mathbf a^{\mathsf T}H\mathbf b)^2\geq
 (\mathbf a^{\mathsf T}H\mathbf a)
 (\mathbf b^{\mathsf T}H\mathbf b).
\]
Continuity extends this to positive real weights and vectors.

Suppose $H$ has two positive eigenvalues, and fix $\mathbf a>0$.
Since $H$ is nonzero and nonnegative, $\mathbf a^{\mathsf T}H\mathbf a>0$.
A two-dimensional positive subspace meets the hyperplane
$\mathbf a^{\mathsf T}H\mathbf z=0$
in a nonzero vector $\mathbf z$ with $\mathbf z^{\mathsf T}H\mathbf z>0$.
For sufficiently small nonzero $\varepsilon$, the vector
$\mathbf b=\mathbf a+\varepsilon\mathbf z$ is positive. Then
\[
 \begin{aligned}
 \mathbf a^{\mathsf T}H\mathbf b
 &=\mathbf a^{\mathsf T}H\mathbf a,\\
 \mathbf b^{\mathsf T}H\mathbf b
 &=\mathbf a^{\mathsf T}H\mathbf a
   +\varepsilon^2\mathbf z^{\mathsf T}H\mathbf z,
 \end{aligned}
\]
contradicting the inequality. Pinning color-$c$ positions handles
higher degrees.
\end{proof}

Every word is counted, and the color projection may have parallel elements.

\section[fixedPalette]{Keeping the number of colors fixed}

We can also remove weights without adding colors. The construction
replaces a binary-list position by bases of an ordinary matroid.

Let $e$ have list $\{a,b\}$. Let $N$ be a rank-$r$ ordinary
matroid on $Z\sqcup\{*\}$, where $*$ is neither a loop nor a coloop,
and put $n=|Z|$. Replace $e$ by $Z$. Above a word with $u_e=a$, use
any basis of $N$ avoiding $*$ as the color-$a$ set in $Z$. Above a
word with $u_e=b$, use $B\setminus\{*\}$ for any basis $B$ containing
$*$. All other positions in $Z$ have color $b$.

\begin{proposition}[Matroid substitution at a binary position]
These expansions form a colored matroid using the same colors.
The quotas of $a$ and $b$ increase by $r-1$ and $n-r$,
respectively. The old entry is recovered from whether $Z$ contains
$r$ or $r-1$ entries of color $a$.
\end{proposition}

\begin{proof}
Associate each expanded word with its original word and its basis
of $N$. Suppose a prescribed new disagreement is $a\to b$ at $z\in Z$.
Symmetric basis exchange through $z$ gives an element $t$ in the other
local basis for which both exchanged sets are bases.

If $t\in Z$, switching $z,t$ is an $a,b$ two-cycle, leaving the
original words fixed. If $t=*$, the original entries at $e$ are $a,b$. Take an
original cyclic exchange through $e$ and replace its arrow by the arrow at
$z$. Both resulting words correspond to an exchanged original word
and an exchanged basis of $N$.

For a prescribed disagreement outside $Z$, take an original cyclic exchange
there. If it uses $e$, symmetric basis exchange through $*$ supplies
a position of $Z$ to replace that arrow. Otherwise the original cycle
works unchanged. Reversing the parents handles a disagreement $b\to a$.
\end{proof}

Suppose that neither $a$ nor $b$ is the distinguished color. To realize
positive integer weights $(s,t)$ at $e$, take $N=U_{t,s+t}$. The two
numbers of local expansions are
\[
 \begin{gathered}
 \binom{s+t-1}{t}=sC,\\
 \binom{s+t-1}{t-1}=tC,\\
 C=\frac{(s+t-1)!}{s!\,t!}>0.
 \end{gathered}
\]
The unweighted completion polynomial after substitution is therefore
$C$ times the original polynomial with weights $(s,t)$ at $e$.

\subsection[quadraticCompatibility]{Pairwise compatibility}

The following lemma combines restrictions to pairs of colors.

\begin{lemma}[Compatibility of quadratic forms]
Let $H_1,\dotsc,H_m$ be nonzero symmetric matrices with nonnegative
entries. Suppose every nonnegative combination of every pair has at
most one positive eigenvalue. Then every nonnegative combination of
all the matrices has at most one positive eigenvalue.
\end{lemma}

\begin{proof}
Suppose $H=\sum_i\alpha_iH_i$ has two positive eigenvalues, with
$\alpha_i\geq0$ and $\sum_i\alpha_i=1$. Choose $\mathbf v>0$ and
$\mathbf y$ with $\mathbf v^{\mathsf T}H\mathbf y=0$ and
$\mathbf y^{\mathsf T}H\mathbf y>0$, as in the preceding proof.

Over the simplex $\beta_i\geq0$, $\sum_i\beta_i=1$, impose the linear
constraint $\sum_i\beta_i\mathbf v^{\mathsf T}H_i\mathbf y=0$.
Maximize $\sum_i\beta_i\mathbf y^{\mathsf T}H_i\mathbf y$.
This feasible compact polytope has an extreme maximizer with at most
two nonzero coordinates: three active coordinates would admit a
nonzero perturbation preserving both linear equalities.
The maximum is positive since $\boldsymbol\alpha$ is feasible.
For $K=\sum_i\beta_iH_i$ at this maximizer, we have
\[
 \begin{gathered}
 \mathbf v^{\mathsf T}K\mathbf v>0,\qquad
 \mathbf v^{\mathsf T}K\mathbf y=0,\\
 \mathbf y^{\mathsf T}K\mathbf y>0.
 \end{gathered}
\]
Thus $K$ is positive definite on the span of $\mathbf v,\mathbf y$,
contradicting the pairwise hypothesis.
\end{proof}

\begin{theorem}[Reduction with a fixed number of colors]
Fix $k\geq2$. The weighted Lorentzian conjecture for all colored
matroids with at most $k$ colors is equivalent to unweighted
Lorentzianity for all binary-list colored matroids with at most $k$
colors and distinguished quota two.
\end{theorem}

\begin{proof}
The forward implication is immediate. For the converse, first restore
all weights for binary-list families of distinguished quota two.
At a position allowing two other colors, matroid substitution realizes
any positive rational weight ratio without changing the color-$c$
variables or adding colors. A common factor in those two weights contributes
the same positive scalar to every word.
At a position with list $\{c,a\}$, extract $w_{e,a}$ and replace
$x_e$ by $x_e/w_{e,a}$. A forced assignment of a different color
contributes a scalar. Approximation and closedness handle real weights.

We next remove the binary-list condition by induction on
$\sum_e\max(|L_e|-2,0)$. At a position allowing at least three colors,
group the terms of the quadratic completion polynomial by their entry there.
Restricting that position to any one or two of its colors gives
a colored matroid with smaller induction parameter. All nonnegative
combinations of the two corresponding sums have at most one positive
Hessian eigenvalue: positive coefficients are implemented by assignment
weights, or by scaling $x_e$ for the color-$c$ sum;
zero coefficients follow by limits. The compatibility lemma applies
to the nonzero Hessians and proves the same condition for
their sum. Its support is already matroidal.

Finally, pinning color-$c$ positions reduces higher degrees to these
quadratic families without adding colors.
\end{proof}
